\section{Results and discussion}
\subsection{RQ1: A narrower observation changes both detection and correctness}
The development camera-ladder diagnostic keeps 7,238 ground-truth buildings in 43 ROIs fixed while changing the rendered footprint. ChangeOS localization recall rises from 22.09\% at cruise to 34.06\% at intermediate and 56.52\% at floor. When unmatched ground-truth buildings are counted as incorrect, binary accuracy rises from 20.01\% to 31.06\% and 53.65\%, respectively. Among matched buildings alone, accuracy is 90.56\%, 91.20\%, and 94.92\%, with macro-F1 of 0.8714, 0.8743, and 0.9232. Matched-only Brier score falls from 0.0657 to 0.0384 between cruise and floor; these probability metrics are not defined over unmatched ground-truth buildings. The large gap between all-GT and matched-only accuracy shows that localization and observation coverage, not just binary classification, dominate the full-building endpoint. Because matched populations change by view, their classification percentages must not replace the fixed-denominator measure.

For the paired cruise-to-floor comparison, 2,696 building outcomes change from incorrect to correct and 261 change from correct to incorrect, a net gain of 2,435 on this development population. The same nominally closer view therefore supplies useful detail and harmful changes. These are binary ChangeOS diagnostics; the archived four-class results in \ref{app:historical} belong to a different backend and are not pooled with them. More importantly, building-level gains do not by themselves imply better task answers.

\subsection{RQ2 and RQ3: Frozen online policies and task answers}
Table~\ref{tab:final-policies} reports the complete online comparison. T1\_TASK has the highest observed accuracy, 75.00\%, compared with 72.29\% for A0\_HOLD. It uses 161 executed reobservations across 480 question-repeat episodes, compared with 419 for random, 923 for always, and 626 for raw entropy. The higher observed score and lower realized reobservation count are useful descriptive results. They are not a matched-cost comparison: policies have the same maximum cap but choose different numbers of actions.

\begin{table}[htbp]
\centering\small
\caption{Frozen final-set online policies. Each row contains the same 160 questions under three generation repeats (480 question-repeat episodes); repeats are not independent questions. Wall time is measured software episode runtime, not flight duration. A3 and A4 are supplemental comparisons.}
\label{tab:final-policies}
\resizebox{\linewidth}{!}{\begin{tabular}{lrrrr}
\toprule
Policy & Correct / 480 & Accuracy (\%) & Reobservations & Wall time (s/episode)\\
\midrule
A0\_HOLD & 347 & 72.29 & 0 & 17.03\\
A1\_RANDOM & 326 & 67.92 & 419 & 33.65\\
A2\_ALWAYS & 307 & 63.96 & 923 & 51.07\\
A3U\_RAW\_ENTROPY & 339 & 70.63 & 626 & 39.61\\
A3\_ENTROPY & 341 & 71.04 & 628 & 38.94\\
A4\_CONFORMAL & 338 & 70.42 & 628 & 38.95\\
T1\_TASK & 360 & 75.00 & 161 & 17.92\\
\bottomrule
\end{tabular}}
\end{table}

The four prespecified paired contrasts are in Table~\ref{tab:final-contrasts}. T1\_TASK exceeds A0\_HOLD by an observed 2.71 percentage points, but the multiplicity-adjusted ROI-cluster interval includes zero. The appropriate claim is that T1 is higher \emph{in this final run}, while a stable advantage over holding remains uncertain. Its contrast with raw entropy is likewise uncertain. The interval against random touches zero at its lower endpoint, so it does not establish a strict positive difference under the declared interval rule. T1 is above always reobserving within the prespecified cluster analysis, but this comparison alone does not prove that its individual movements caused the gain.

\begin{table}[htbp]
\centering\small
\caption{Prespecified T1\_TASK paired accuracy contrasts. Differences are percentage points; 98.75\% ROI-cluster percentile intervals adjust for the four primary comparisons. All questions and repeats from a sampled ROI remain together.}
\label{tab:final-contrasts}
\begin{tabular}{lrr}
\toprule
Contrast & Difference (pp) & 98.75\% interval (pp)\\
\midrule
T1 minus A0\_HOLD & $+2.71$ & $[-2.04,\,+7.50]$\\
T1 minus A1\_RANDOM & $+7.08$ & $[0.00,\,+14.70]$\\
T1 minus A2\_ALWAYS & $+11.04$ & $[+0.86,\,+21.38]$\\
T1 minus A3U\_RAW\_ENTROPY & $+4.38$ & $[-4.33,\,+12.96]$\\
\bottomrule
\end{tabular}
\end{table}

Additional observations can alter the answer in either direction. Within their own episode histories, T1\_TASK records 36 corrected and 21 harmed answers; A2\_ALWAYS records 74 corrected and 132 harmed. These transition labels are not paired counterfactual effects versus A0\_HOLD. They nevertheless show why an always-reobserve policy can spend the largest action budget while finishing with the lowest accuracy. A1\_RANDOM executes 419 reobservations and reaches 67.92\%; A3U\_RAW\_ENTROPY executes 626 and reaches 70.63\%. Neither action count nor classifier uncertainty alone orders final task quality.

The episode traces show concrete instances in both directions. In the first generation repeat, a Hurricane Michael damage question (ID ending \texttt{damage\_0001\_9430}) changes from an initially incorrect ``no damage'' answer to the correct ``damaged'' answer after two T1 reobservations. A Hurricane Michael spatial question (ID ending \texttt{spatial\_0092\_5452}) changes from the correct ``north'' to the incorrect ``southwest'' after one. These examples describe recorded answer sequences, not isolated causal effects of the movement.

\subsection{RQ4 and RQ5: Task dependence, bottlenecks, and scope of the gain}
Task-type outcomes reveal uneven effects. For T1\_TASK versus A0\_HOLD, damage accuracy is 75.76\% versus 66.67\% and count accuracy is 78.29\% versus 75.97\%. Presence is identical at 93.18\%; spatial accuracy is lower for T1, 41.38\% versus 43.68\%. These are descriptive, correlated subgroup summaries, not four additional adjusted tests. They are consistent with the design premise that a wide context can already suffice for presence while a more detailed target view may help damage, but the spatial result limits any universal benefit claim.

\begin{table}[htbp]
\centering\small
\caption{Task evidence requirements and descriptive final accuracy. A0 and T1 use the same questions and three repeats. The evidence requirements follow the task definitions; they are not experimentally isolated causes of the accuracy differences.}
\label{tab:task-evidence}
\begin{tabularx}{\linewidth}{@{}lXrr@{}}
\toprule
Task & Evidence requirement & A0 (\%) & T1 (\%) \\
\midrule
Presence & Regional coverage and positive evidence & 93.18 & 93.18 \\
Damage & Target identity and local damage detail & 66.67 & 75.76 \\
Count & ROI coverage and consistent object aggregation & 75.97 & 78.29 \\
Spatial & Reference geometry and nearest-target identity & 43.68 & 41.38 \\
\bottomrule
\end{tabularx}
\end{table}

Table~\ref{tab:task-evidence} connects these outcomes to benchmark construction. Narrowing a view allocates more source pixels to a building, but it can remove other buildings needed for a count or the context needed to identify the nearest damaged target. The benchmark therefore exposes competing requirements for detail and coverage rather than ranking every narrower view as better. The platform's joint record of ROI coverage, predicted objects, actions, and answers makes such trade-offs inspectable; it does not by itself resolve their causal contributions.

Event outcomes are also heterogeneous. T1\_TASK versus A0\_HOLD is 54.86\% versus 58.33\% in Hurricane Michael, 84.31\% versus 77.12\% in the Moore tornado, and 83.06\% versus 79.23\% in Nepal flooding. The pooled numerical gain is not reproduced in every event, and three events cannot support an unseen-disaster generalization claim. Accuracy across T1's three generation repeats is 75.00\%, 74.38\%, and 75.63\%; repeat stability does not remove ROI or event uncertainty.

Development diagnostics provide limited bottleneck evidence without forming an additive oracle decomposition. A deterministic answerer using predicted structured evidence reaches 66.87\% on an earlier development run. Full-ROI ground-truth evidence with a deterministic rule answerer reconstructs 160/160 development reference answers, checking the question/label pipeline rather than online performance. Their source fingerprints and interventions differ, and current-visible GT with Qwen, full-task GT with Qwen, and a full-coverage always-floor reference remain unmeasured. The final T1/A0 comparison therefore isolates neither the perception error nor the language-generation error as a unique cause.

The observed T1 operating point combines a modest numerical accuracy gain over holding with fewer executed reobservations than other active policies. Its uncertainty interval, event variation, and incomplete oracle chain favor a platform-and-mechanism interpretation rather than a demonstrated generally superior task-conditioned algorithm. A randomized common-budget sweep would be needed to estimate a genuine accuracy--cost curve; the seven realized operating points cannot supply one by rebinning episodes after execution.
